% rexxformat.tex -- combined listings setup for NetRexx and rxas
%
% This file deliberately defines languages and styles separately.  Use
% style=netrexxcode or style=rxascode on \lstinputlisting / lstlisting.
%
% It also avoids calling \fontspec directly from \lstset; a named font
% command plus \selectfont is more reliable with listings and OpenType fonts.

\ifdefined\codefont
\else
  \newfontfamily\codefont{IBM Plex Mono}
\fi

\lstdefinelanguage{NetRexx}
{morekeywords={abstract,adapter,binary,case,catch,class,constant,dependent,deprecated,digits,do,else,end,engineering,extends,final,finally,for,forever,if,implements,indirect,import,implements,indirect,inheritable,interface,iterate,label,leave,loop,method,native,nop,numeric,options,otherwise,over,package,parent,parse,private,properties,protect,public,return,returns,rexx,say,scientific,set,digits,form,select,shared,signal,signals,sourceline,static,super,then,this,until,used,upper,volatile,when,where,while},
sensitive=false,
extendedchars=true,
morecomment=[s]={/*}{*/},
morecomment=[l]{--},
morecomment=[s]{/**}{*/},
morestring=[b]",
morestring=[d]",
morestring=[b]',
morestring=[d]'}

\lstdefinelanguage{rxas}
{morekeywords={
appendchar,
addf,
addi,
addi,
amap,
amap,
and,
append,
addf,
bcf,
bcf,
bct,
bct,
bctnm,
bctnm,
beq,
bge,
bge,
bgt,
bgt,
ble,
ble,
blt,
blt,
bne,
bne,
bpoff,
bpon,
br,
brf,
brt,
brtpandt,
beq,
brtpt,
brtf,
call,
cnop,
call,
concchar,
concat,
concat,
concat,
call,
copy,
divf,
dec,
dec0,
dcall,
dec2,
dec1,
divf,
divi,
divi,
dropchar,
divf,
exit,
exit,
exit,
fadd,
fadd,
fcopy,
fdiv,
fdiv,
fdiv,
feq,
fformat,
fgt,
fgt,
fgt,
fgte,
fgte,
fgte,
flt,
flt,
flt,
flte,
flte,
flte,
fmult,
fmult,
fndblnk,
fndnblnk,
fne,
fne,
fpow,
fpow,
fsex,
fsub,
fsub,
fsub,
ftob,
ftoi,
ftos,
feq,
getstrpos,
gmap,
gettp,
gmap,
getbyte,
hexchar,
iadd,
iadd,
igt,
iand,
icopy,
idiv,
idiv,
idiv,
ieq,
ieq,
igt,
igte,
iand,
igte,
igte,
ilt,
ilt,
ilt,
ilte,
ilte,
ilte,
imod,
imod,
imod,
imult,
imult,
inc,
inc0,
inc1,
ine,
igt,
ine,
inot,
inot,
ior,
ior,
ipow,
ipow,
ipow,
irand,
irand,
isex,
ishl,
ishl,
ishr,
ishr,
isub,
isub,
isub,
itof,
itos,
ixor,
ixor,
inc2,
linkattr,
load,
linkattr,
loadsettp,
load,
load,
load,
link,
loadsettp,
loadsettp,
metaloadedmodules,
map,
map,
metadecodeinst,
metaloadcalleraddr,
metaloaddata,
metaloadedprocs,
metaloadfoperand,
metaloadinst,
metaloadioperand,
metaloadmodule,
metaloadpoperand,
metaloadsoperand,
move,
mtime,
multf,
multi,
metalinkpreg,
multi,
multf,
nsmap,
nsmap,
nsmap,
null,
nsmap,
not,
or,
padstr,
pmap,
pmap,
poschar,
ret,
ret,
ret,
ret,
ret,
ret,
rseq,
rseq,
readline,
sappend,
say,
say,
say,
say,
sayx,
sayx,
sconcat,
sconcat,
sconcat,
scopy,
seq,
seq,
setortp,
setstrpos,
settp,
sgt,
sgt,
sgt,
sgte,
sgte,
sgte,
slt,
slt,
slt,
slte,
slte,
slte,
sne,
sne,
stof,
stoi,
strchar,
strchar,
strlen,
strlower,
strupper,
subf,
subf,
subf,
subi,
subi,
substcut,
substr,
substring,
swap,
say,
time,
transchar,
triml,
triml,
trimr,
trimr,
trunc,
trunc,
unlink,
unmap,
xtime,
  },
sensitive=false,
extendedchars=true,
morecomment=[s]={/*}{*/},
morecomment=[l]{*},
morecomment=[s]{/**}{*/},
morestring=[b]",
morestring=[d]",
morestring=[b]',
morestring=[d]'}

% Common code styling.  Keep the font command in every style that listings may
% apply while tokenising; otherwise delimiters and comments can be positioned
% or rendered differently from the surrounding text.
\lstdefinestyle{rexxcommon}{
  captionpos=none,
  tabsize=3,
  keywordstyle=\codefont\color{nrorange},
  commentstyle=\codefont\color{nrgrey},
  stringstyle=\codefont\color{nrgreen},
  numberstyle=\tiny,
  numbersep=5pt,
  breaklines=true,
  showstringspaces=false,
  index=[1][keywords],
  columns=fullflexible,
  keepspaces=true,
  upquote=true,
  emph={label}
}

\lstdefinestyle{netrexxcode}{
  style=rexxcommon,
  language=NetRexx,
  alsolanguage=rexx,
  numbers=none,
  basicstyle=\codefont\fontsize{10}{12}\selectfont
}

\lstdefinestyle{rxascode}{
  style=rexxcommon,
  language=rxas,
  numbers=left,
  basicstyle=\codefont\fontsize{10}{12}\selectfont
}

% Optional default.  Comment this out if each listing selects its own style.
\lstset{
style=netrexxcode,
 literate={*}{{{\codefont\char"002A}}}1
 }
